\subsection{Environmental impacts}

The environmental impacts of \acp{bss} have been a central part of their justification~\cite{Shaheen2010, Fishman2016, Midgley2011}. These impacts include transport emissions, energy use, material consumption, and end-of-life waste and recycling. While early academic discussions often focused on \ac{ghg} reductions from substituting car trips, recent research has broadened the perspective to include the full life cycle of shared bicycles, docking infrastructure, batteries, and logistic operations.

        \subsubsection{Reduction of transport emissions derived from modal substitution} 
        \label{subsubsection:emissions_reduction}

        Transport emissions encompass two related but distinct pollutant categories: \acp{ghg} such as carbon dioxide (\ce{CO2}), methane (\ce{CH4}), and nitrous oxide (\ce{N2O}), which accumulate in the atmosphere, and local air pollutants such as nitrogen oxides (\ce{NOx}), hydrocarbon (\ce{HC}), particulate matter (\ce{PM_{2.5}} and \ce{PM_{10}}), and carbon monoxide (\ce{CO}), which are short-lived and primarily affect urban air quality. Both categories are produced by motorized transport modes, but their implications differ: \ac{ghg} reductions contribute to climate mitigation, while pollutant reductions directly improve street-level air quality.

        Early environmental assessments often reported large climate benefits, typically by estimating the total distance travelled by bike-sharing and then attributing these kilometres to alternative modes using assumed substitution patterns. Emission savings were calculated by applying mode-specific emission factors to the replaced kilometres and subtracting the negligible emissions of cycling.

        In Barcelona, for example, \textit{Bicing} was estimated to avoid more than 9,000 tonnes of \ce{CO2} in 2011, but assuming that 90\% of users were new cyclists and that all of their trips would otherwise have been made by car~\cite{Rojas-Rueda2011}. Later work in Shanghai estimated savings of 25,000 tonnes of \ce{CO2} in 2016. However, the study assumed that all trips shorter than 1~km replaced walking, while longer trips replaced taxi, and the annual city-wide results were extrapolated from a single month of data by adjusting for the sample ratio (56.62\%), the company’s market share (56.56\%), and scaling up to twelve months~\cite{Zhang2018}. At the same time, Zhang~\cite{Zhang2018} acknowledged that the environmental benefits of bike-sharing are inherently localized in time and space, since they are directly correlated with patterns of bicycle use rather than being uniform across the city or year. Other studies reported even larger figures. Qiu and He~\cite{Qiu2018}, for instance, estimated a reduction of 616,036 tonnes of \ce{CO2} in 2020 compared to the 2015 value, assuming that 75\% of bike-sharing kilometres were replacing car trips. An analysis of New York’s \textit{Citi Bike} marked some progress in this field by mapping trips onto the street network rather than using straight-line distances, but its estimates of 30,070 tonnes of \ce{CO2} reductions between 2014 and 2017 still rested on deterministic distance-based thresholds borrowed from other research~\cite{Chen2022c}.

        Taken together, these approaches tend to overstate benefits mainly because they make strong assumptions about mode substitution (more details on modal substitution in Section~\ref{sec:impacts_of_bss}). To address these shortcomings, studies refined how modal substitution is modelled. Kou et al.~\cite{Kou2020} developed a trip-level probabilistic model that combined observed bike-share distances with city-specific travel survey data and life-cycle emission factors, assigning each trip to a most likely replaced mode. Using this method, they estimated annual reductions across eight U.S. cities in 2016 ranging from 41 to 5,417 tonnes of \ce{CO2}, or roughly 0.3--0.6~kg per trip. Similarly, Saltykova~\cite{Saltykova2022}, analyzing dockless bike-sharing in Chengdu, used traveller survey responses to quantify modal substitution. They found that 39\% of bike-share trips replaced bus journeys, 13.5\% replaced subway, and only 10.9\% replaced car trips. With these figures, the avoided emissions fell markedly compared to a “cars-only” counterfactual—demonstrating how omitting transit substitution inflates climate benefits by a factor of two or more.

        A more refined approach was proposed by Li et al.~\cite{Liu2021}, who developed a high-resolution discrete-choice framework matching over 23 million dockless bike-share trips in Shanghai to realistic modal alternatives using travel time, distance, and cost information from online routing services. Their results showed that most trips substituted for walking (30.7\%) or public transport (50\%), with car replacement below 2\%, leading to estimated annual savings of 116,000 tonnes of \ce{CO2}. This progression—from simple distance rules to survey-informed probabilities, and finally to discrete-choice models—illustrates a gradual shift toward more behaviourally grounded and context-sensitive estimates of the environmental benefits.

        Research has also estimated reductions in local air pollutants following the same logic as with \ac{ghg} emissions, although such assessments are far less common. For example, in Shanghai, bike-sharing was estimated to cut \ce{NOx} emissions by 64 tonnes~\cite{Zhang2018}; while in Beijing, \ce{SO2}, \ce{CO}, and \ce{NO2} emissions were reduced by 22.5, 58.64, and 1,586.7 tonnes, respectively, and \ce{PM_{2.5}} and \ce{PM_{10}} by 10.3 tonnes~\cite{Qiu2018}.


        \subsubsection{Life-cycle and material impacts}

        \acp{bss} do not only generate impacts through avoiding transport emissions. One early study that highlighted this was conducted by Fishman, Washington, and Haworth~\cite{Fishman2014}, who found that London’s bike-sharing scheme in 2012 produced higher overall motor vehicle usage because of the rebalancing operations required to ensure the service. Therefore, to account for the whole impact of a \ac{bss} more factors should be accounted. The manufacturing of bicycles, docks, and batteries, the maintenance and logistics operations, and the handling of waste when the hardware is no longer useful, all require material- and energy-intensive processes, which contribute to embodied emissions.

        To capture these impacts, studies typically employ Life-Cycle Assessments (LCA), usually based on the ISO 14040 and 14044 standards~\cite{LCA2006a, LCA2006b}. These standards define how to account for the burdens associated with manufacturing, operation, maintenance, and end-of-life management (Figure~\ref{fig:LCA_bss}). Importantly, full LCAs should also incorporate the benefits from modal substitution, so that avoided emissions from displaced car or public transport trips are evaluated alongside the embodied burdens of system operation.
        
        \begin{figure}[htbp!]
            \centering
            \includegraphics[width=1\textwidth, trim={0cm 0.3cm 0cm 1cm}, clip]{0_plots/figure_LCA.pdf}
            \caption{Aspects that should be considered in a \ac{bss} LCA. Figure based on Luo et al.~\cite{Luo2019}, adapted and elaborated by the author.}
            \label{fig:LCA_bss}
        \end{figure}

        Within the LCA framework, the life-cycle inventory phase quantifies all relevant material and energy flows entering and leaving the system. These inventories constitute the empirical foundation for the subsequent impact assessment stage, where the quantified flows are translated into environmental indicators. To perform this translation, researchers apply established methods such as ReCiPe 2008~\cite{Bonilla-Alicea2020}, ReCiPe 2016~\cite{Felipe-Falgas2022}, or GREET~\cite{Cazzola2020, Krauss2022}. These approaches allow impacts to be expressed either at the midpoint level, which reports specific categories such as climate change or particulate matter formation with relatively low uncertainty, or at the endpoint level, which aggregates them into broader areas of protection such as human health, ecosystems, and resources, with higher uncertainty but greater communicability. Earlier studies of shared micromobility often limited their scope to \ac{ghg} emissions, while more recent LCAs studies present both midpoint and endpoint indicators~\cite{Sun2022, Chen2024, Calan2024}.

        By applying these methodologies, researchers can compute comparable indicators across systems. At the midpoint level, results are typically reported per service delivered, in grams of \ce{CO2}-eq per passenger-kilometer (pkm) of bike travel~\cite{Luo2019, Chen2024, Bonilla-Alicea2020}. The term \ce{CO2}-eq (“carbon dioxide equivalent”) expresses the combined impact of different \acp{ghg}, such as \ce{CO2}, \ce{CH4}, and \ce{N2O}, by converting them into the amount of \ce{CO2} that would cause the same warming effect, based on their global warming potentials. This metric enables comparison across systems, but also with other transport modes, since it accounts for how intensively bicycles are used~\cite{Sun2022, Krauss2022}. Reported values for complete \acp{bss} LCAs range from roughly 30 to 150~g~\ce{CO2}-eq/pkm~\cite{Zhu2023, Felipe-Falgas2022, Cazzola2020, Sun2022, Calan2024}. This wide variation reflects differences in system type, usage and operational practices, and whether all life-cycle phases are included in the analysis.

        The first step of an LCA is assessing the \textbf{manufacturing phase}, which covers the extraction of raw materials and production of bicycles, and docking stations and batteries where applicable. This is modelled by compiling life-cycle inventories of material inputs—such as aluminium and steel for frames, lithium and cobalt for batteries, plastics for saddles and tyres, and electronics for smart bikes or stations—and linking them with emission factors from databases like Ecoinvent~\cite{Felipe-Falgas2022} or CPCD~\cite{Zhu2023}. 

        Manufacturing often emerges as the largest single contributor, accounting for around 50\% of total life-cycle emissions in station-based systems~\cite{Calan2024, Sun2022}. When comparing station-based and dockless systems with \acp{e-b}, the production of the hardware accounts 61\% of the life-cycle emissions for station-based systems and 52\% for dockless systems~\cite{Luo2019}. Moreover, Bonilla-Alicea~\cite{Bonilla-Alicea2020} quantified manufacturing burdens and found that dockless systems' bicycles with on-board electronics had approximately 3.7 times the climate change impact of docked bicycles, and about 2.7 times their total endpoint impacts. The increase was largely driven by the electronic subsystems (solar panels, printed circuit boards, communication modules, and batteries), highlighting that these components are candidates for technological improvement in terms of environmental impact. On the other hand, shared e-bikes add further burdens from battery and motor production, with batteries alone contributing 20–30\% of embodied \acp{ghg}~\cite{Zhu2023}. Material choices also play a decisive role. Earlier LCAs of private bicycles showed that aluminium frames carried higher production impacts~\cite{Dave2010, Leunberger2010, Luna2016} than alternatives such as steel or bamboo~\cite{Agyekum2017}, a finding echoed in shared fleets where heavier, vandal-resistant frames and integrated electronics further raise embodied emissions.

        In electrically assisted systems, these manufacturing burdens become even more significant, since batteries require periodic replacement and their production is a major hotspot. Evidence from related micromobility fleets shows the scale of these effects: in shared e-scooters, lithium batteries account for 10–19\% of manufacturing emissions~\cite{Calan2024}, and battery packs may need to be replaced up to 2.4 times during an e-bike’s service life~\cite{Calan2024}. 
        
        The \textbf{use phase} of \acp{bss} is not limited to the act of cycling itself, which is practically emission-free, but also involves several supporting operations. LCAs typically account for rebalancing, battery charging, maintenance, and modal substitution savings. 
        
        Among these, \textbf{rebalancing} consistently emerges as the dominant contributor, responsible for 36\% of total life-cycle emissions in station-based systems and up to 73\% in dockless ones~\cite{Luo2019}, which corresponded to with average values around 118~g~\ce{CO2}-eq/pkm for dockless systems and 65--70~g~\ce{CO2}-eq/pkm for station-based ones~\cite{Luo2019}.

        The high burden of rebalancing arises because it requires motorized service vehicles to redistribute bicycles across the city. Its impact depends on several factors: the number of bicycles in the system, overall system size, the types of vehicles used for redistribution, and how often rebalancing occurs~\cite{Bonilla-Alicea2020}. To quantify this burden, LCAs model the distance travelled by rebalancing vehicles relative to the service delivered, often expressed as service-vehicle kilometers per bicycle-kilometer. Estimates range from as low as 0.012 km of van travel per bike-km in station-based systems~\cite{Felipe-Falgas2022} to 0.10–0.26 km in dockless fleets~\cite{Sun2022, Cazzola2020}. Some systems, for instance, rely on conventional vans with intensities around 300~g~\ce{CO2}-eq/km~\cite{Chen2020LifeCycle, Hollingsworth2019}, while others have begun experimenting with electric vans or cargo bikes for battery collection and charging, significantly reducing the associated emissions~\cite{Felipe-Falgas2022, Calan2024}. More broadly, improving rebalancing efficiency has become a research field of its own, with numerous studies developing optimization models and operational strategies to minimize environmental and financial costs.

        Moreover, \textbf{battery charging} should also be considered. As Sun et al.~\cite{Sun2022} emphasize, electrically assisted systems incur burdens not only from the collection and redistribution of bicycles, but also from the energy required to recharge batteries. The direct electricity demand of shared e-bikes is relatively modest—around 0.014--0.017~kWh per km, or roughly 1.5~kWh per 100~km~\cite{Sun2022}—yet its impact depends strongly on the carbon intensity of the grid, ranging from about 10--20~g~\ce{CO2}-eq/km in low-carbon contexts to more than 70~g~\ce{CO2}-eq/km in coal-dominated regions~\cite{Zhu2023, Calan2024}. Although small compared to the emissions from rebalancing logistics, charging remains an integral part of the use-phase footprint and interacts with operational choices such as battery swapping, charging frequency, and the type of service vehicles used.

        Another component of the use phase are \textbf{maintenance operations}, which include routine activities such as component replacement, pumping, lubricating, and cleaning~\cite{Zhu2023}. Because operator records are rarely available~\cite{Calan2024}, LCAs typically rely on interviews~\cite{Zhu2023}, technical manuals, or web sources~\cite{Bonilla-Alicea2020} to estimate service frequencies and part lifetimes. As a result, maintenance is often underrepresented in system assessments. When explicitly included, maintenance is usually found to contribute around 10–15\% of total life-cycle \acp{ghg} in conventional station-based systems~\cite{Sun2022, Chen2024, deBortoli2021Environmental}. These shares can rise in dockless fleets, where higher rates of wear and vandalism accelerate the need for component replacement~\cite{Luo2019}.

        Overall, the relative importance of maintenance depends strongly on assumptions about vehicle lifetimes and component durability. Shorter lifespans for bicycles and batteries increase the share of maintenance, whereas strategies that extend durability and improve repairability can substantially mitigate impacts~\cite{Zhu2023, Calan2024}.      

        Beyond operational burdens, the use phase in most LCAs also credits \textbf{avoided emissions from modal substitution}. In practice, this means assigning each bike-share trip a counterfactual mode (e.g., car, public transport, walking) and subtracting the corresponding life-cycle emissions that would have been produced otherwise. Recent LCAs operationalize this with survey-based mode shares or trip data: in Washington, D.C., Chen et al.~\cite{Chen2024} incorporate mode-shift information to compute a net balance where avoided emissions are weighed against manufacturing and operations. In Barcelona, Felipe-Falgas et al.~\cite{Felipe-Falgas2022} fuse an LCA inventory with self-reported substitution and show that shared e-bikes can increase acp{ghg} when they primarily replace public transport or walking rather than cars. At a broader scale, Sun et al.~\cite{Sun2022} formalize the accounting as the difference between the bike-share emissions and a weighted mix of displaced modes (plus new trips), reporting results in g~\ce{CO2}-eq/pkm.

        Several syntheses reinforce this integrated “burdens minus benefits” framing. Krauss et al.~\cite{Krauss2022} adapted mode-shift survey data and LCA factors to compute the net sustainability impact of shared micromobility across six cities, showing that outcomes depend critically on the share of car and ride-hailing trips displaced versus low-carbon modes. Methodological guidance from the ITF/OECD likewise treats modal shift as essential to comparability, recommending city-specific substitution profiles and transparent functional units (pkm) when reporting net outcomes~\cite{Cazzola2020}. Together, these studies demonstrate that whether bike sharing reduces \acp{ghg} in the use phase depends less on cycling’s direct emissions (which are near zero) and more on what trips it replaces and how operations are managed.

        The final step of an LCA is assessing the \textbf{end-of-life phase}, which covers the collection of retired bicycles, batteries, docks, and electronic components by operators, followed by recycling and disposal processes. In Krauss et al.~\cite{Krauss2022}, this phase is included but represents only a minor share of total life-cycle emissions, since most impacts stem from vehicle manufacturing and operational services. Similar to the maintenance phase, it is often underrepresented in LCAs~\cite{Calan2024}. To address this gap, both Krauss et al.~\cite{Krauss2022} and Sun et al.~\cite{Sun2023} assume that 80–90\% of metal components in shared micromobility fleets can be recycled, with the remainder incinerated or landfilled.

        Several studies have examined how different treatment pathways affect environmental outcomes. Zhu~\cite{Zhu2023}, for example, conducted a LCA of shared e-bikes in China and compared two disposal scenarios. Over the full three-year life cycle of a single vehicle, the net \ac{ghg} reduction was about 487~kg~\ce{CO2}-eq if non-recyclable parts were incinerated, but only 434~kg~\ce{CO2}-eq if they were landfilled, where some materials could persist for decades. These values represent net balances that integrate all life-cycle stages, including production, operation, and substitution effects, not disposal alone. In Washington, Chen~\cite{Chen2024} similarly modelled both incineration and landfill pathways, but reported stage-specific results: disposal contributed 42.5~kg~\ce{CO2}-eq and 48.6~kg~\ce{CO2}-eq per bicycle, respectively. By contrast, manufacturing accounted for about 63\% of the life-cycle total, producing roughly five times more emissions than disposal. Despite different reporting approaches, both studies indicate that incineration performs better than landfill, and that disposal plays only a minor role compared to production.

        Beyond metals and frames, batteries remain a particular hotspot: their production is already a major contributor, and their retirement is complicated by hazardous materials and recycling inefficiencies~\cite{Felipe-Falgas2022, Calan2024}. Without robust recycling schemes, battery disposal risks offsetting part of the emission savings from modal substitution.

        Empirical evidence highlights large variability in recycling performance. Reported recovery rates range from 30\% to 90\%, depending on the material and regulatory context~\cite{Sun2022, Calan2024}. Metals such as aluminium and steel are typically assumed to have high recovery potential, while plastics, electronics, and batteries are more difficult to recycle and are frequently treated as waste. In Washington, D.C., Chen~\cite{Chen2024} showed that moving from landfill to waste-to-energy reduced the end-of-life footprint of a shared bike by about 6~kg~\ce{CO2}-eq, while systematic recycling of metals offered larger offsets.

        In practice, fleet attrition and oversupply can further amplify end-of-life burdens. When 20\% of bikes are lost annually to vandalism or theft, overall fleet emissions may increase by a quarter because additional units must be manufactured to replace prematurely retired vehicles~\cite{Calan2024}. The most striking illustration is the emergence of “bike-share graveyards” in Chinese cities during the 2017–2018 dockless boom~\cite{Sun2022}, when uncontrolled fleet expansion left millions of unused bicycles piled up in open lots (Fig.~\ref{fig:bss_graveyard}). In these cases, low utilization meant that manufacturing emissions could not be amortised across sufficient trips, while weak recycling systems left much of the material unrecovered~\cite{Luo2019, Sun2022}.

        Overall, the end-of-life phase tends to contribute only a small fraction of total \acp{ghg}. Its importance lies less in the direct emissions than in the opportunity to close material loops, recover critical resources, and avoid the accumulation of hazardous waste. From a sustainability perspective, robust recycling infrastructures, modular product design, and right-sized fleets are key strategies to improve end-of-life outcomes~\cite{Krauss2022, Calan2024}.

        \begin{figure}    
            \centering
            \includegraphics[width=0.8\textwidth]{0_plots/graveyard.jpg}
            \caption{Bike-sharing graveyard in Xiamen, China. Source: https://www.amusingplanet.com/2018/05/chinas-bicycle-graveyards.html}
            \label{fig:bss_graveyard}
        \end{figure}